\section{Results}

The results are organised around the main analysis outputs: preprocessing and
tumour-sample retention, neighbourhood consensus across feature--distance
representations, patient-referenced and pan-cancer similarity networks,
marker-gene structure, functional enrichment, tumour-to-cell-line ranking, and
cell-line-to-tumour similarity ranking. The final subsection reports the
cell-line-to-tumour similarity ranking, in
which 54 of 56 biological cell-line groups had a same-lineage tumour as the
top-ranked sample.

\subsection{Preprocessing retained high-purity solid tumour cohorts and reduced data-origin separation}

\subsubsection{Tumour purity filtering retained solid tumour cohorts}
At the predefined tumour purity threshold of $\geq 0.7$, 914 of 1{,}231 breast cancer samples, 244 of 282 neuroblastoma samples, and 66 of 66 retinoblastoma samples were retained (Fig.~\ref{fig:brca_purity_analysis}; \suppFigsRef{fig:nbl_purity_analysis}{fig:mpn_purity_analysis}). In contrast, no samples from the Burkitt lymphoma dataset or the TCGA diffuse large B cell lymphoma cohort were retained. In the breast cancer cohort, tumour purity was inversely associated with immune and stromal scores (Spearman correlation \(r = -0.90\) and \(r \approx -0.85\), respectively; both \(p < 2.2 \times 10^{-16}\); Fig.~\ref{fig:brca_purity_analysis}). Equivalent ESTIMATE-based purity checks were available for neuroblastoma and retinoblastoma: neuroblastoma showed inverse immune and stromal correlations (Spearman \(r = -0.95\) and \(r = -0.91\), respectively; both \(p < 2.2 \times 10^{-16}\)), and retinoblastoma showed weaker but significant inverse immune and stromal correlations (Spearman \(r = -0.44\), \(p = 2.4 \times 10^{-4}\), and \(r = -0.45\), \(p = 1.6 \times 10^{-4}\), respectively; \suppFigsRef{fig:nbl_purity_analysis}{fig:mpn_purity_analysis}).


\begin{figure}[H]
\centering
\begin{subfigure}{0.48\linewidth}
  \includegraphics[width=\linewidth]{Thesis_writeup/Figures/tumour_purity_analysis/breast_cancer_tpa/brca_sample_retention.pdf}
  \caption{}
\end{subfigure}
\hfill
\begin{subfigure}{0.48\linewidth}
  \includegraphics[width=\linewidth]{Thesis_writeup/Figures/tumour_purity_analysis/breast_cancer_tpa/brca_tumour_distribution.pdf}
  \caption{}
\end{subfigure}
\vspace{0.4cm}
\begin{subfigure}{0.48\linewidth}
  \includegraphics[width=\linewidth]{Thesis_writeup/Figures/tumour_purity_analysis/breast_cancer_tpa/brca_tumour_immune_score.pdf}
  \caption{}
\end{subfigure}
\hfill
\begin{subfigure}{0.48\linewidth}
  \includegraphics[width=\linewidth]{Thesis_writeup/Figures/tumour_purity_analysis/breast_cancer_tpa/brca_tumour_purity_stromalScore.pdf}
  \caption{}
\end{subfigure}
\caption{
Tumour purity analysis in TCGA breast cancer. 
(a) Sample retention before and after tumour purity filtering.
(b) Distribution of tumour purity estimates; the dashed line indicates the applied threshold.
(c) Tumour purity versus ESTIMATE immune score.
(d) Tumour purity versus ESTIMATE stromal score.
}
\label{fig:brca_purity_analysis}
\end{figure}

\subsubsection{UMAP visualisation showed reduced separation by data origin after preprocessing}

In UMAP space, raw transcriptomic profiles separated primarily by data
origin, with cell lines forming clusters distinct from patient tumour samples
(Fig.~\ref{fig:umap_combined_vertical}a). After preprocessing, profiles showed
less visual separation by data origin, while retained solid tumour samples
separated mainly by tumour lineage (Fig.~\ref{fig:umap_combined_vertical}b).
The curated 257-gene pan-cancer feature set, defined from recurrent and
isolate-associated marker genes, showed a similar lineage-separated structure
(Fig.~\ref{fig:umap_combined_vertical}c).

\begin{figure}[H]
\centering
\raggedright \textbf{(a)}\par
\centering
\includegraphics[width=0.7\linewidth]{Thesis_writeup/Figures/umap_raw_cell_tumour_60616.pdf}

\vspace{1em}
\raggedright \textbf{(b)}\par
\centering
\includegraphics[width=0.7\linewidth]{Thesis_writeup/Figures/umap_batch_corrected_cell_tumour_60616.pdf}

\vspace{1em}
\raggedright \textbf{(c)}\par
\centering
\includegraphics[width=0.7\linewidth]{Thesis_writeup/Figures/umap_feature_set_breast_cancer_257.pdf}
\caption{
UMAP projections of transcriptomic structure before and after preprocessing.
(a) Raw expression matrix.
(b) Processed expression matrix.
(c) Curated 257-gene feature set.
}
\label{fig:umap_combined_vertical}
\end{figure}

\subsection{Neighbourhood consensus differed across feature--distance representations and lineages}

Neighbourhood consensus varied across feature--distance representations in all
three tumour lineages. Breast cancer was evaluated with ten representations,
including PAM50, whereas neuroblastoma and retinoblastoma were evaluated with
the nine shared unsupervised representations. The reported outcome was the
fraction of evaluated cell-line--tumour neighbourhood relationships supported
at \(p_{\mathrm{consensus}} \geq 0.7\).

In the breast cancer cohort, the fraction of consensus-supported neighbourhood relationships ranged from 17\% for WGCNA total connectivity (\(kTotal\)) under correlation distance to 57\% for the MAD--Euclidean representation (Fig.~\ref{fig:feature_distance_combined}a). Dispersion-based feature sets paired with Euclidean distance yielded the highest fractions of consensus-supported neighbour assignments, including MAD--Euclidean at 57\%, mean absolute deviation--Euclidean at 52\%, and entropy--Euclidean at 49\%. The \(kTotal\)--correlation representation had the lowest consensus fraction.

In the neuroblastoma cohort, the fraction of consensus-supported neighbourhood relationships ranged from 59\% for the entropy--correlation representation to 92\% for the \(kTotal\)--correlation representation (Fig.~\ref{fig:feature_distance_combined}b). Connectivity-based and rank-based representations had the highest support, including \(kTotal\)--correlation at 92\% and Spearman-connectivity--correlation at 91\%. Entropy representations had the lowest consensus fractions, and correlation distance exceeded Euclidean distance across most feature classes.

In the retinoblastoma cohort, the fraction of consensus-supported neighbourhood relationships ranged from 18\% for the Spearman-connectivity--correlation representation to 93\% for the mean absolute deviation--Euclidean representation (Fig.~\ref{fig:feature_distance_combined}c). Euclidean representations showed higher consensus-supported neighbourhood fractions than correlation variants across feature classes, with the highest support from mean absolute deviation--Euclidean at 93\%, entropy--Euclidean at 85\%, variance--Euclidean at 82\%, and MAD--Euclidean at 80\%.

\begin{figure}[H]
\centering
\begin{minipage}{0.9\textwidth}
    \centering
    \raggedright \textbf{(a)}\par
    \centering
    \includegraphics[width=0.88\linewidth]{Thesis_writeup/Figures/breast_cancer/feature_dis_breast_cancer.pdf}

    \vspace{1em}
    \raggedright \textbf{(b)}\par
    \centering
    \includegraphics[width=0.88\linewidth]{Thesis_writeup/Figures/neuroblastoma/feature_dis_neuroblastoma.pdf}

    \vspace{1em}
    \raggedright \textbf{(c)}\par
    \centering
    \includegraphics[width=0.88\linewidth]{Thesis_writeup/Figures/retinoblastoma/feature_dis_retinoblastoma.pdf}
\end{minipage}
\caption{
Comparison of feature selection methods for neighbourhood analysis from cell lines to patient tumour samples across three tumour lineages.
The x axis shows compact feature--distance labels in the format \texttt{feature\_distance}, where \texttt{euc} denotes Euclidean distance and \texttt{corr} denotes Pearson correlation distance. Bar height indicates the fraction of pairs linking cell lines to patient tumour samples with neighbourhood consensus $p_{\text{consensus}} \geq 0.7$. Colours distinguish feature selection methods, and distance suffixes distinguish the distance metric used for patient sample ranking.
(a) Breast cancer cohort ($n = 818$ patient tumour samples; 29 breast cancer cell lines).
(b) Neuroblastoma cohort ($n = 244$ patient tumour samples; 18 cell lines).
(c) Retinoblastoma cohort ($n = 66$ patient tumour samples; 9 biological cell-line groups).
Feature selection abbreviations: Variance, variance-ranked genes; HVG, highly variable genes; MAD, median absolute deviation-ranked genes; MeanAbsDev, mean absolute deviation-ranked genes; Entropy, Shannon entropy-ranked genes; PCA, genes ranked by principal component analysis loading score; Spearman, genes ranked by mean absolute Spearman connectivity; MX, connectivity--variance score combining mean absolute Spearman connectivity with scaled gene expression variance; kTotal, WGCNA total network connectivity-ranked genes; PAM50, breast cancer PAM50 subtype gene set (panel a only).
}
\label{fig:feature_distance_combined}
\end{figure}

\subsection{Maximum consensus feature--distance representations varied across cell lines}

Maximum consensus feature--distance representations varied across cell lines
within each tumour lineage. In the breast cancer cohort, the maximum fraction
of consensus-supported neighbour pairs per cell line ranged from 0.45 for
CAL\_51 to 1.00 (Fig.~\ref{fig:per_cellline_feature_distance_combined}a).
Seven breast cancer cell lines reached perfect consensus. The MAD--Euclidean
representation was the most frequent maximum consensus representation,
occurring in eleven cell lines.

In the neuroblastoma cohort, per cell line maximum consensus values ranged
from 0.65 for SH\_SY5Y to 1.00 (Fig.~\ref{fig:per_cellline_feature_distance_combined}b).
Eleven of eighteen cell lines reached perfect consensus. The
Spearman-connectivity--correlation representation was the most frequent
maximum consensus representation, occurring in fourteen of eighteen cell lines.

In the retinoblastoma cohort, maximum consensus assignments were more
concentrated than in the other two lineages (Fig.~\ref{fig:per_cellline_feature_distance_combined}c).
The maximum fraction of consensus-supported neighbour pairs ranged from 0.71
for RBL\_14 to 1.00. Seven of nine biological cell-line groups reached perfect consensus. The
mean absolute deviation--Euclidean representation was the most frequent
maximum consensus representation, occurring in eight of nine biological cell-line groups.
Y\_79 reached perfect consensus under seven Euclidean representations.

\begin{figure}[H]
\centering
\begin{minipage}{0.9\textwidth}
    \centering
    \raggedright \textbf{(a)}\par
    \centering
    \includegraphics[width=0.88\linewidth]{Thesis_writeup/Figures/breast_cancer/per_cellline_feature_dis_breast_cancer.pdf}

    \vspace{1em}
    \raggedright \textbf{(b)}\par
    \centering
    \includegraphics[width=0.88\linewidth]{Thesis_writeup/Figures/neuroblastoma/per_cellline_feature_dis_neuroblastoma.pdf}

    \vspace{1em}
    \raggedright \textbf{(c)}\par
    \centering
    \includegraphics[width=0.88\linewidth]{Thesis_writeup/Figures/retinoblastoma/per_cellline_feature_dis_retinoblastoma.pdf}
\end{minipage}
\caption{
Per cell line maximum consensus feature--distance representation across three tumour lineages.
Each row corresponds to one cell line. Points mark the feature--distance representation that maximised the fraction of consensus-supported neighbourhood pairs for that cell line. Compact representation labels follow the \texttt{feature\_distance} format, where \texttt{euc} denotes Euclidean distance and \texttt{corr} denotes Pearson correlation distance. Point colour denotes the feature selection method, point shape denotes the distance metric, and values on the right axis show the maximum fraction of neighbourhood pairs with $p_{\text{consensus}} \geq 0.7$.
(a) Breast cancer cohort ($n = 818$ patient tumour samples; 29 breast cancer cell lines).
(b) Neuroblastoma cohort ($n = 244$ patient tumour samples; 18 cell lines).
(c) Retinoblastoma cohort ($n = 66$ patient tumour samples; 9 biological cell-line groups).
Feature selection abbreviations are as defined in Fig.~\ref{fig:feature_distance_combined}.
}
\label{fig:per_cellline_feature_distance_combined}
\end{figure}

\subsection{Patient-referenced cell-line similarity networks}

Patient-referenced support-threshold networks showed lineage-specific
cell-line neighbourhood structures across breast cancer, neuroblastoma, and
retinoblastoma. The cohort-specific networks differed in component structure,
isolated profiles, and high-betweenness nodes.

\subsubsection{Breast cancer cell line similarity network}

The breast cancer support-threshold consensus network contained 29 breast cancer cell lines
partitioned into two connected components (Fig.~\ref{fig:brca_consensus_dsmz}).
The bridge-like annotations marked BT\_474, EFM\_192C, and ETCC\_006 as
breast cancer anchor cell lines.

\begin{figure}[H]
\centering
\includegraphics[width=\linewidth]{Thesis_writeup/Figures/breast_cancer/Fig_dsmz_brca_consensus_graph.pdf}
\caption{
Support-threshold consensus network for breast cancer cell lines.
The network comprises two connected components.
}
\label{fig:brca_consensus_dsmz}
\end{figure}

\subsubsection{Neuroblastoma cell line similarity network}

The neuroblastoma support-threshold consensus network included all 18 analysed
cell lines; 15 retained support-threshold edges and formed two connected
components, while LS, SH\_SY5Y, and SK\_N\_BE\_2 appeared as isolates
(Fig.~\ref{fig:nbl_consensus_dsmz_graph}). The dominant component contained 13
cell lines, and a smaller component linked GI\_ME\_N and
CCLF\_PEDS\_0051\_T. SIMA, KELLY, and CHP\_126 had the highest unnormalised
betweenness centrality within the dominant component.

\begin{figure}[H]
\centering
\includegraphics[width=\linewidth]{Thesis_writeup/Figures/neuroblastoma/Fig_nbl_consensus_dsmz_graph.pdf}
\caption{
Support-threshold consensus network for neuroblastoma cell lines.
The displayed network shows the 15 edge-retaining cell lines: a 13-node
component encompassing the majority of neuroblastoma cell lines and a two-node component
linking GI\_ME\_N and CCLF\_PEDS\_0051\_T. LS, SH\_SY5Y, and SK\_N\_BE\_2 were
isolates and are not displayed in this figure.
}
\label{fig:nbl_consensus_dsmz_graph}
\end{figure}

\subsubsection{Retinoblastoma cell line similarity network}

The retinoblastoma support-threshold consensus network was analysed at
biological-group level after arithmetic mean aggregation of replicate
tumour-neighbourhood consensus vectors (Fig.~\ref{fig:rbl_consensus_dsmz_graph}).
RBL\_20\_1 and RBL\_20\_3 were represented as RBL\_20, and RBL\_15\_2 and
RBL\_15\_4 were represented as RBL\_15. The resulting support-threshold
network contained nine biological retinoblastoma cell-line groups and four
retained edges.

\begin{figure}[H]
\centering
\includegraphics[width=\linewidth]{Thesis_writeup/Figures/retinoblastoma/Fig_rbl_dsmz_consensus_graph.pdf}
\caption{
Support-threshold consensus network for retinoblastoma biological cell-line
groups. Replicate profile identifiers for RBL\_15 and RBL\_20 were collapsed
before graph construction by arithmetic mean aggregation of their
tumour-neighbourhood consensus vectors. The network comprises nine biological
cell-line groups and four retained support-threshold edges.
}
\label{fig:rbl_consensus_dsmz_graph}
\end{figure}

\subsection{Resolved neighbour graphs contained cell line neighbourhood components and isolates}

Graph-based resolution produced resolved neighbour graphs that retained the
analysed cell-line node set, including isolates and nodes without retained
support-threshold edges. For the final retinoblastoma graph, replicate profile
identifiers were collapsed to biological cell-line groups before graph
construction. The resulting components, isolates, and component anchors are
reported below.

\subsubsection{Resolved breast cancer neighbourhood components}

The resolved graph comprised four multi-node connected components of sizes 14,
6, 3, and 3, together with three isolated cell lines
(Fig.~\ref{fig:brca_resolved_components}).
Component anchors and isolated cell lines are labelled in the figure; the
isolated cell lines were DU\_4475, IPH\_926, and SK\_BR\_3.

\begin{figure}[H]
\centering
\includegraphics[width=\linewidth]{Thesis_writeup/Figures/breast_cancer/breast_cancer_resolved_components.pdf}
\caption{
Resolved components for the graph-based breast cancer neighbourhood network.
The graph comprises four multi-node components of sizes 14, 6, 3, and 3,
together with three isolated cell lines.
}
\label{fig:brca_resolved_components}
\end{figure}

\subsubsection{Resolved neuroblastoma neighbourhood components}

In the resolved neuroblastoma neighbour graph, the neuroblastoma cell lines were
partitioned into an 11-node component and a two-node component containing
MHH\_NB\_11 and LAN\_5 (Fig.~\ref{fig:nbl_resolved_components}). The
GI\_ME\_N--CCLF\_PEDS\_0051\_T pair was retained. LS, SH\_SY5Y, and SK\_N\_BE\_2
appeared as isolates.

\begin{figure}[H]
\centering
\includegraphics[width=\linewidth]{Thesis_writeup/Figures/neuroblastoma/nbl_resolved_components.pdf}
\caption{
Resolved components for the graph-based neuroblastoma neighbourhood network.
The graph comprises an 11-node component, two two-node components, and three
isolated cell lines.
}
\label{fig:nbl_resolved_components}
\end{figure}

\subsubsection{Resolved retinoblastoma neighbourhood components}

In the resolved retinoblastoma neighbour graph, the retinoblastoma cell lines
were analysed as nine biological groups after arithmetic mean aggregation of
replicate profile vectors (Fig.~\ref{fig:rbl_resolved_components}). The graph
contained one five-node component comprising RBL\_13, RBL\_15, RBL\_18, RBL\_20,
and RBL\_30, with RBL\_18 selected as both the most-connected node and
bridge-like anchor. RBL\_7, WERI\_RB1, Y\_79, and RBL\_14 appeared as isolates.

\begin{figure}[H]
\centering
\includegraphics[width=\linewidth]{Thesis_writeup/Figures/retinoblastoma/rbl_resolved_components.pdf}
\caption{
Resolved components for the graph-based retinoblastoma neighbourhood network.
The resolved graph was constructed at biological-group level after arithmetic
mean aggregation of replicate profile tumour-neighbourhood consensus vectors.
It comprises a five-node component (RBL\_13, RBL\_15, RBL\_18, RBL\_20,
RBL\_30) and four isolated biological cell-line groups (RBL\_7, WERI\_RB1,
Y\_79, RBL\_14).
}
\label{fig:rbl_resolved_components}
\end{figure}

\subsection{Differential expression analysis of resolved components and isolates}

Marker gene structure was then characterised across the resolved components
and isolates.

The final differential expression analysis included 21 contrasts: eight in breast cancer, seven in neuroblastoma, and six in retinoblastoma. Marker yield differed by contrast type and lineage.

Isolate contrasts retained between 20 and 50 markers in breast cancer, between
6 and 50 markers in neuroblastoma, and between 21 and 50 markers in
retinoblastoma. Anchor contrasts retained between 177 and 200 markers in breast
cancer, between 17 and 200 markers in neuroblastoma, and between 82 and 160
markers in retinoblastoma (\suppTableRef{tab:dge_contrasts}).

Within each disease, recurrence filtering retained qualifying markers that
appeared in at least two contrasts. Recurrence filtering reduced the breast
cancer marker set from 1{,}007 genes observed in at least one contrast to 87
genes in the strict recurrent core, the neuroblastoma marker set from 582 genes
to 20 genes, and the retinoblastoma marker set from 385 genes to one gene.

Recurrence filtering and isolate-associated marker rescue produced the final
257-gene pan-cancer feature set, containing all 108 genes in the strict
recurrent core (87 breast cancer, 20 neuroblastoma, one retinoblastoma). This
feature set was used for pan-cancer network construction and similarity ranking.
The complete gene lists are provided in
\suppTablesRef{tab:dge_recurrent_genes}{tab:pancancer_257_genes}.


\subsection{Pan-cancer transcriptomic similarity network contained five Louvain communities with partial neuroblastoma--retinoblastoma assignment overlap}

Using the final 257-gene pan-cancer feature set, the pan-cancer transcriptomic
similarity network contained 167 cell-line RNA-seq profiles and 2{,}130
nearest-neighbour Spearman similarity edges, spanning breast cancer (BRCA),
haematological comparator profiles (HEME), neuroblastoma (NBL), and
retinoblastoma (RBL). The network formed a single connected component
(Fig.~\ref{fig:network}).

\subsubsection{Community structure, topology, and lineage purity}

Louvain community detection partitioned the network into five communities
ranging from 18 to 56~nodes (Table~\ref{tab:summary}).
The network had modularity $Q = 0.629$, and lineage assortativity was
$r = 0.801$.
Within-community edge weights and degree measures are reported
in Table~\ref{tab:summary}.

Four of five communities had purity $\geq 0.879$ (Table~\ref{tab:purity}). Haematological comparator profiles (HEME) were the majority lineage in
communities~2, 3, and~5, each with purity of 1.000. The lineage-coloured
network layout is shown in Fig.~\ref{fig:network}. Community~4 contained the
main cross-lineage assignment pattern, with neuroblastoma (NBL) as the majority lineage.
In this separate pan-cancer profile-level network, community~4 contained 11
retinoblastoma (RBL) RNA-seq profiles corresponding to nine unique
retinoblastoma cell lines, including replicate libraries for RBL-15 and
RBL-20.

\subsubsection{Lineage-discordant cell line profiles}

Of the 167~nodes with lineage annotation, 15~nodes (9.0\%),
representing 12~unique cell lines, were assigned to a community whose majority
lineage differed from their own annotation (Table~\ref{tab:discordant}).
Most discordant assignments followed two patterns: HT-1080, GI\_ME\_N, and
CCLF\_PEDS\_0051\_T were assigned to the breast-cancer (BRCA)-majority
community, whereas retinoblastoma (RBL) profiles were assigned to the
NBL-majority community.

\begin{figure}[H]
  \centering
  \includegraphics[width=0.95\textwidth]{Thesis_writeup/Figures/Fig_pan_cancer_cell_line_similarity_network_lineage_community}
  \caption{\textbf{Pan-cancer transcriptomic similarity network of 167 cell line RNA-seq profiles.} Nodes are coloured by cancer lineage: breast cancer (BRCA), orange;
    haematological comparator profiles (HEME), sky blue; neuroblastoma (NBL), bluish green; retinoblastoma (RBL), reddish purple; unknown, grey.
    Node size is proportional to degree centrality. Edges connect nearest-neighbour RNA-seq profiles in Spearman correlation space; the layout places more strongly connected profiles closer together.
    Network properties: 167~nodes; 2{,}130~edges; five Louvain
    communities; modularity $Q = 0.629$; lineage assortativity $r = 0.801$.}
  \label{fig:network}
\end{figure}

\begin{table}[H]
\centering
\caption{\textbf{Structural properties of the five communities in the 167-node
  RNA-seq profile similarity network.} Within edges denote edges with
  both endpoints in the same community. Mean edge weight denotes the mean
  Spearman correlation of within community edges.}
\label{tab:summary}
\begin{tabular}{crrrrrr}
\toprule
Community & $n$ nodes & Within edges & Mean edge weight &
  Mean degree & Median degree & Max degree \\
\midrule
1 & 33 & 365 & 0.737 & 23.2 & 21 & 32 \\
2 & 33 & 375 & 0.780 & 26.7 & 25 & 42 \\
3 & 56 & 691 & 0.735 & 28.3 & 24 & 52 \\
4 & 27 & 306 & 0.690 & 23.3 & 24 & 27 \\
5 & 18 & 110 & 0.752 & 22.3 & 21 & 32 \\
\bottomrule
\end{tabular}
\end{table}

\begin{table}[H]
\centering
\caption{\textbf{Community purity by majority lineage.} Purity is defined as
  the majority lineage count divided by community size. All 167~nodes in
  this community table carried lineage labels.}
\label{tab:purity}
\begin{tabular}{cclrr}
\toprule
Community & $n$ nodes & Majority lineage & Majority count & Purity \\
\midrule
1 & 33 & BRCA & 29 & 0.879 \\
2 & 33 & HEME & 33 & 1.000 \\
3 & 56 & HEME & 56 & 1.000 \\
4 & 27 & NBL  & 16 & 0.593 \\
5 & 18 & HEME & 18 & 1.000 \\
\bottomrule
\end{tabular}
\end{table}

\begin{table}[H]
\centering
\caption{\textbf{Lineage-discordant cell line profiles ($n = 12$ unique cell lines;
  $n = 15$ nodes).} Each entry carries a lineage label that differs from
  the majority lineage of its assigned community. Where a cell line contributed
  more than one RNA-seq profile, the profile count is noted.}
\label{tab:discordant}
\begin{tabular}{clllr}
\toprule
Community & Cell line & Lineage & Majority lineage & $n$ nodes \\
\midrule
1 & HT-1080           & HEME & BRCA & 2 \\
1 & GI\_ME\_N         & NBL & BRCA & 1 \\
1 & CCLF-PEDS-0051-T  & NBL & BRCA & 1 \\
4 & RBL-7             & RBL & NBL  & 1 \\
4 & RBL-13            & RBL & NBL  & 1 \\
4 & RBL-14            & RBL & NBL  & 1 \\
4 & RBL-15            & RBL & NBL  & 2 \\
4 & RBL-18            & RBL & NBL  & 1 \\
4 & RBL-20            & RBL & NBL  & 2 \\
4 & RBL-30            & RBL & NBL  & 1 \\
4 & WERI-RB1          & RBL & NBL  & 1 \\
4 & Y-79              & RBL & NBL  & 1 \\
\bottomrule
\end{tabular}
\end{table}


\subsection{Functional enrichment of marker gene sets}

Functional enrichment of disease-specific and shared marker gene sets
returned lineage-associated biological terms. Component and contrast
identities for the tested marker sets are listed in
\suppTablesRef{tab:supp_component_annotations}{tab:dge_contrasts}; selected
enriched terms are shown in Fig.~\ref{fig:enrichment_top_terms_heatmap}, and
complete term-level enrichment outputs are listed in the Supplementary data
files.

\subsubsection{Breast cancer}

Downregulated marker genes from ETCC\_006 relative to component~5 were enriched for cell periphery
(GO:CC; $p = 1.18 \times 10^{-16}$, $n = 92$) and plasma membrane
(GO:CC; $p = 4.62 \times 10^{-14}$, $n = 85$). Marker genes from
EFM\_192C relative to component~0 had similar
enrichment for cell periphery (GO:CC; $p = 5.78 \times 10^{-17}$,
$n = 89$) and plasma membrane (GO:CC; $p = 1.99 \times 10^{-14}$,
$n = 82$).

HCC\_1937 isolate markers were enriched for 2q37 copy number variation
syndrome in the all marker set ($p = 1.39 \times 10^{-16}$, $n = 10$)
and downregulated set ($p = 6.91 \times 10^{-17}$, $n = 10$), together
with ETS1 TF motif enrichment (TF; \(p = 2.01 \times 10^{-3}\),
\(n = 19\)). The BRCA
disease recurrent downregulated set was enriched for tube development
(GO:BP; $p = 6.70 \times 10^{-12}$, $n = 25$) and multicellular
organismal process (GO:BP; $p = 1.45 \times 10^{-11}$, $n = 54$).

\subsubsection{Neuroblastoma}

Upregulated marker genes from CCLF\_PEDS\_0051\_T relative to component~0 were enriched for Interleukin 10
signalling (Reactome; $p = 1.76 \times 10^{-18}$, $n = 12$), SARS-CoV-2
signalling network map (WikiPathways; $p = 7.48 \times 10^{-16}$,
$n = 19$), defense response (GO:BP; $p = 1.64 \times 10^{-14}$,
$n = 39$), cytokine-cytokine receptor interaction (KEGG;
$p = 2.19 \times 10^{-14}$, $n = 16$), and inflammatory response
(GO:BP; $p = 3.44 \times 10^{-14}$, $n = 28$).

Within component~1, downregulated markers from CHP\_126 were enriched for p300 TF motif
(TF; $p = 1.66 \times 10^{-10}$, $n = 87$) and CPBP TF motif
(TF; $p = 4.38 \times 10^{-10}$, $n = 92$). Within component~1, IMR\_32 marker genes were enriched for multicellular organism development
(GO:BP; $p = 5.33 \times 10^{-5}$, $n = 34$) and anatomical structure
development (GO:BP; $p = 6.01 \times 10^{-5}$, $n = 38$).

LS downregulated markers were enriched for ZNF148 TF motif (TF;
$p = 2.72 \times 10^{-7}$, $n = 21$), neuron development
(GO:BP; $p = 3.47 \times 10^{-6}$, $n = 15$), and neuron differentiation
(GO:BP; $p = 4.57 \times 10^{-6}$, $n = 16$). SH\_SY5Y returned no
significant enrichment in the all-marker input set.

\subsubsection{Retinoblastoma}

The WERI\_RB1 all-marker isolate set was enriched for regulation of multicellular
organismal process (GO:BP; $p = 5.35 \times 10^{-5}$, $n = 22$),
signalling (GO:BP; $p = 1.06 \times 10^{-4}$, $n = 31$), and cell
communication (GO:BP; $p = 1.20 \times 10^{-4}$, $n = 31$). The
WERI\_RB1 upregulated set also returned death receptor activity
(GO:MF; $p = 1.70 \times 10^{-4}$, $n = 3$) and SMAD5 TF motif
(TF; $p = 3.41 \times 10^{-4}$, $n = 16$).

RBL\_14 upregulated markers were enriched for basement membrane collagen
network terms, including network forming collagen trimer, collagen type IV
trimer, and collagenous component of basement membrane
(GO:CC; $p = 1.19 \times 10^{-4}$). The Y\_79 downregulated set returned
retinal photoreceptor expression annotations (HPA; Retina, photoreceptor
cells $\geq$ Low; $p = 5.48 \times 10^{-4}$, $n = 4$). The RBL\_18 anchor
contrast for the main biological-group component returned CORUM complex
enrichment, including the PXN--ITGB5--PTK2 complex. Selected retinoblastoma
enrichment terms are included in Fig.~\ref{fig:enrichment_top_terms_heatmap}.

\subsubsection{Pan-cancer shared gene-set enrichment}

The final 257-gene pan-cancer feature set was enriched for T cell modulation
and desmoplasia in pancreatic cancer
(WikiPathways; $p = 2.46 \times 10^{-6}$,
$n = 8$), epithelial to mesenchymal transition in colorectal cancer
(WikiPathways; $p = 3.22 \times 10^{-4}$, $n = 10$), and focal adhesion
(WikiPathways; $p = 3.56 \times 10^{-4}$, $n = 11$). The TF source also
returned a broad source-level annotation (TF; $p = 2.39 \times 10^{-27}$,
$n = 237$) and broad TF motif enrichments, including Spi-B
(TF; $p = 1.94 \times 10^{-27}$, $n = 237$), TR4
(TF; $p = 8.23 \times 10^{-27}$, $n = 230$), Kid3
(TF; $p = 1.03 \times 10^{-26}$, $n = 236$), and CPBP
(TF; $p = 2.94 \times 10^{-25}$, $n = 231$).

Within the union of the disease-level strict recurrent cores, the
downregulated set contained 98 genes and was enriched for developmental terms,
including multicellular organism development
(GO:BP; \(p = 2.52 \times 10^{-14}\), \(n = 52\)), tube development
(GO:BP; \(p = 2.81 \times 10^{-14}\), \(n = 28\)), cell migration
(GO:BP; \(p = 9.39 \times 10^{-13}\), \(n = 30\)), anatomical structure
morphogenesis (GO:BP; \(p = 5.65 \times 10^{-13}\), \(n = 39\)), and
E2F-2 TF motif (TF; \(p = 7.82 \times 10^{-18}\), \(n = 91\)). The
all-marker set contained 108 genes and returned related developmental
terms. The upregulated set contained 19 genes and returned eight significant
terms, including hypertrophic cardiomyopathy
(KEGG; \(p = 1.94 \times 10^{-3}\), \(n = 3\)),
dilated cardiomyopathy
(KEGG; \(p = 2.19 \times 10^{-3}\), \(n = 3\)), and
TFII-I TF motif (TF; $p = 4.54 \times 10^{-3}$, $n = 13$). The
isolate-rescued, non-recurrent subset returned IL13RA1--TRAF3IP1 CORUM complex
enrichment (CORUM; $p = 1.66 \times 10^{-2}$, $n = 2$). Selected enriched
terms across disease-specific and shared gene-set categories are shown in
Fig.~\ref{fig:enrichment_top_terms_heatmap}.

\begin{figure}[H]
    \centering
    \includegraphics[width=\textwidth]{Thesis_writeup/Figures/Fig_enrichment_top_terms_heatmap.pdf}
    \caption{
    Selected enriched terms across disease-specific and pan-cancer marker gene
    sets. Rows show representative enriched terms retained using predefined priority
    terms and enrichment-score ranking from the complete enrichment results.
    Columns correspond to disease-specific per-contrast marker sets,
    within-disease recurrent marker sets, recurrent core gene sets
    (down-regulated, up-regulated, and combined directions),
    the final 257-gene pan-cancer feature set, and the isolate-rescued subset.
    Tile intensity represents capped \( -\log_{10} \) g:SCS-corrected enrichment
    \(p\)-value. Empty cells indicate that the term was not retained for display
    in that category after term selection; they do not necessarily indicate that
    the term was tested and non-significant. TF motif and CORUM complex
    enrichments are reported in the complete enrichment output files but are not
    necessarily retained in this reduced display.
    }
    \label{fig:enrichment_top_terms_heatmap}
\end{figure}


\subsection{Tumour-to-cell-line transcriptomic similarity ranking}

In the same curated pan-cancer feature space, 58 cell line models were ranked
for each patient tumour sample across the breast cancer, neuroblastoma, and
retinoblastoma cohorts
(Fig.~\ref{fig:ecdf_model_prioritisation_combined}). The lowest-median-rank
lineage-matched models ranked ahead of the non-focal lineage distribution
across all three cohorts.

In the breast cancer cohort ($n = 818$), the lowest-median-rank breast cancer models shown in Fig.~\ref{fig:ecdf_model_prioritisation_combined} had median ranks between 8 and 9 and top 10 frequencies between 448/818 and 508/818 samples. KPL-1 had the lowest median rank among the selected breast cancer models.

In the neuroblastoma cohort ($n = 244$), the lowest-median-rank neuroblastoma models shown in Fig.~\ref{fig:ecdf_model_prioritisation_combined} had median ranks between 2 and 6 and top 10 frequencies between 211/244 and 225/244 samples. NBLS had the lowest median rank, whereas IMR32 had the highest top 10 frequency.

In the retinoblastoma cohort ($n = 66$), the lowest-median-rank retinoblastoma models shown in Fig.~\ref{fig:ecdf_model_prioritisation_combined} had median ranks between 2 and 4. RBL-14 ranked within the top 10 for 65/66 samples, whereas RBL-18 and both RBL-20 replicate libraries ranked within the top 10 for all 66 samples in this profile-level ranking analysis.

Across all three cohorts, the ECDF curves for selected lineage-matched models
were left-shifted relative to the median non-focal lineage curve
(Fig.~\ref{fig:ecdf_model_prioritisation_combined}).

\begin{figure}[H]
\centering
\includegraphics[width=\linewidth]{Thesis_writeup/Figures/ecdf_plots/ecdf_model_prioritisation_combined.pdf}
\caption{
Empirical cumulative distributions of cell line similarity ranks across patient tumour samples.
For each patient tumour sample, 58 cell line models were ranked by Spearman similarity within the curated pan-cancer feature space, with rank 1 indicating the highest similarity.
(a) Breast cancer cohort ($n = 818$).
(b) Neuroblastoma cohort ($n = 244$).
(c) Retinoblastoma cohort ($n = 66$).
Coloured step curves show selected lineage-matched models. The grey dashed curve shows the median ECDF of non-focal lineage models, and the grey ribbon shows the interquartile range across non-focal lineage models. The vertical dotted line marks rank 10, and the horizontal dotted line marks the median cumulative fraction. The x axis is shown on a $\log_{10}$ scale and extends to all 58 ranked models.
(d) Fraction of patient tumour samples in which each selected model ranked within the top 10. Points show observed fractions, horizontal bars show 95\% Wilson confidence intervals, and labels show exact counts.
}
\label{fig:ecdf_model_prioritisation_combined}
\end{figure}

\subsection{Cell-line-to-tumour similarity ranking}

In the cell-line-to-tumour direction, each of 56 biological cell-line groups
was used to rank 1{,}128 patient tumour samples. The top-ranked tumour lineage
matched the labelled cell-line lineage in 54 of 56 biological cell-line groups,
corresponding to a top-ranked lineage accuracy of 96.4\%
(95\% Wilson CI 87.9\%--99.0\%) and a balanced accuracy of 96.3\%
(Fig.~\ref{fig:cellline_centred_tumour_class_similarity}a). Breast cancer groups
achieved 29/29 correct top-ranked matches (100\%; 95\% Wilson CI 88.3\%--100\%),
and retinoblastoma groups achieved 9/9 correct matches (100\%; 95\% Wilson CI
70.1\%--100\%). Neuroblastoma groups achieved 16/18 correct matches (88.9\%;
95\% Wilson CI 67.2\%--96.9\%); the two exceptions were LS and
CCLF\_PEDS\_0051\_T, both of which ranked a breast cancer sample first.

The complete ranking output, with every patient tumour sample ranked for
each biological cell-line group, is provided in
\suppTableRef{tab:supp_cellline_tumour_rankings}.

Same-lineage tumour samples had median rank percentiles of 3.0\% for
retinoblastoma groups (IQR 1.5\%--4.4\%), 10.9\% for neuroblastoma groups
(IQR 5.5\%--16.2\%), and 36.3\% for breast cancer groups (IQR
18.2\%--54.4\%) (Fig.~\ref{fig:cellline_centred_tumour_class_similarity}c). At the
top 10\% rank cutoff, the median fraction of same-lineage tumours matched per
group was 100\% for retinoblastoma, 45.9\% for neuroblastoma, and 13.7\% for
breast cancer (Fig.~\ref{fig:cellline_centred_tumour_class_similarity}b,d).

The Spearman correlation score of the top-ranked tumour match ranged from
0.74 to 0.86 for retinoblastoma groups, 0.35 to 0.75 for neuroblastoma
groups, and 0.26 to 0.63 for breast cancer groups
(Fig.~\ref{fig:cellline_centred_confidence_scores}b). Confidence margins
between the top-ranked and subsequent matches varied across groups and
lineages; WERI\_RB1 and Y\_79 had the largest rank~1 minus rank~2 margins
among retinoblastoma groups (Fig.~\ref{fig:cellline_centred_confidence_scores}a).

Top-50 tumour component mapping showed that breast cancer groups were largely
dominated by the large BRCA dominant Component~1, with subsets mapping to
smaller BRCA dominant components
(Fig.~\ref{fig:cellline_centred_component_mapping}). Neuroblastoma and
retinoblastoma groups had component compositions distributed across multiple
components, including smaller lineage-dominant components for subsets of
groups, while Component~1 contributed to many top-50 neighbourhoods.

Replicate-level sensitivity analysis for RBL\_15 and RBL\_20 showed similar
same-lineage rank-percentile patterns across independent libraries from the
same cell line (Fig.~\ref{fig:cellline_centred_replicate_sensitivity}).

\begin{figure}[H]
\centering
\includegraphics[width=\linewidth]{Thesis_writeup/Figures/cellline_centred/Fig_cellline_centred_tumour_class_similarity.pdf}
\caption{
Cell-line-to-tumour similarity ranking analysis.
For each of 56 biological cell-line groups, 1{,}128 patient tumour samples were ranked by Spearman correlation within the curated pan-cancer feature space, with higher correlation indicating a closer match and rank 1 indicating the highest similarity.
(a) Top-1 lineage confusion matrix; cells show counts and row percentages. Overall accuracy 96.4\% (95\% Wilson CI 87.9\%--99.0\%), balanced accuracy 96.3\%.
(b) Mean precision@$k$ across biological cell-line groups, stratified by lineage. Dashed lines indicate lineage prevalence null expectations; ribbons show 95\% bootstrap confidence intervals; $n = 1{,}128$ tumour samples.
(c) Same-lineage tumour rank percentile distributions per biological cell-line group; each point represents one same-lineage tumour sample; lower percentile indicates a closer match.
(d) Lineage composition of the top $k = 50$ tumour neighbourhood per biological cell-line group; groups are ordered by same-lineage precision.
BRCA = breast cancer ($n = 29$ biological cell-line groups, $n = 818$ tumour samples); NBL = neuroblastoma ($n = 18$ biological cell-line groups, $n = 244$ tumour samples); RBL = retinoblastoma ($n = 9$ biological cell-line groups, $n = 66$ tumour samples).
}
\label{fig:cellline_centred_tumour_class_similarity}
\end{figure}

\begin{figure}[H]
\centering
\includegraphics[width=\linewidth]{Thesis_writeup/Figures/cellline_centred/Fig_cellline_centred_confidence_scores.pdf}
\caption{
Confidence margins and top-10 Spearman score distributions for cell-line-to-tumour similarity ranking.
(a) Score difference between the top-ranked tumour match and the second-ranked match (rank~1 minus rank~2; left) and the tenth-ranked match (rank~1 minus rank~10; right), stratified by cell-line lineage. Points show biological cell-line groups; boxes indicate the interquartile range.
(b) Distribution of Spearman correlation scores among the top 10 ranked tumour samples per biological cell-line group; groups are ordered by median top-10 score. Cell-line lineage: BRCA = breast cancer; NBL = neuroblastoma; RBL = retinoblastoma.
}
\label{fig:cellline_centred_confidence_scores}
\end{figure}

\begin{figure}[H]
\centering
\includegraphics[width=\linewidth]{Thesis_writeup/Figures/cellline_centred/Fig_cellline_centred_component_mapping.pdf}
\caption{
Consensus network component composition of the top 50 tumour neighbourhood per biological cell-line group.
For each biological cell-line group, the 50 highest-ranked tumour samples were assigned to consensus network components; stacked bars show the fraction of the 50 tumours belonging to each component. Groups are ordered by same-lineage fraction in the top 50. Component labels indicate the dominant tumour lineage, component size ($n$), and purity.
}
\label{fig:cellline_centred_component_mapping}
\end{figure}

\begin{figure}[H]
\centering
\includegraphics[width=0.6\linewidth]{Thesis_writeup/Figures/cellline_centred/Fig_cellline_centred_replicate_sensitivity.pdf}
\caption{
Replicate library sensitivity before group-level collapse.
RBL-15 and RBL-20 were each sequenced across two independent libraries. Points show the median same-lineage rank percentile per replicate library.
}
\label{fig:cellline_centred_replicate_sensitivity}
\end{figure}
